\section*{Capítulo \ref{ch:g1:day-and-night} --- El día y la noche: el Sol}

\begin{solution}{exo:g1:day-and-night:1}
El \omterm{def:g1:day-and-night:daynight}{amanecer}, el mediodía, la tarde, el \omterm{def:g1:day-and-night:daynight}{atardecer}, la noche.
\end{solution}

\begin{solution}{exo:g1:day-and-night:2}
Al mediodía el Sol está en su punto más alto --- y la \omterm{def:g1:light-and-shadows:shadow}{sombra} es lo más
corta que va a ser.
\end{solution}

\begin{solution}{exo:g1:day-and-night:3}
Siempre por el mismo lado. Para comprobarlo: fíjate mañana por la
mañana en por qué ventana entra la primera \omterm{def:g1:light-and-shadows:source}{luz} del Sol --- será la
misma ventana que hoy.
\end{solution}

\begin{solution}{exo:g1:day-and-night:4}
Falso. El Sol brilla toda la noche --- al otro lado de la Tierra. La
noche llega porque nuestro lado de la Tierra, que gira, se ha apartado
de la \omterm{def:g1:light-and-shadows:source}{luz}, igual que la nuca cuando le das la espalda a la lámpara.
\end{solution}

\begin{solution}{exo:g1:day-and-night:5}
Para ellos es de noche, así que lo más probable es que estén
durmiendo --- su lado de la Tierra está de espaldas al Sol mientras el
nuestro lo mira de frente.
\end{solution}

\begin{solution}{exo:g1:day-and-night:6}
Con la \omterm{def:g1:light-and-shadows:source}{luz} en la mejilla derecha y girando todavía, lo siguiente es
que la \omterm{def:g1:light-and-shadows:source}{luz} le llegue a toda la cara --- el \omterm{def:g1:day-and-night:daynight}{amanecer} (la mañana) para
tu nariz.
\end{solution}

\begin{solution}{exo:g1:day-and-night:7}
El suelo que tenemos bajo los pies. La Tierra gira y nos lleva
consigo; el Sol solo \emph{parece} cruzar nuestro cielo porque somos
nosotros los que giramos.
\end{solution}

\begin{solution}{exo:g1:day-and-night:8}
Siete \omterm{def:g1:day-and-night:daynight}{amaneceres} y siete noches --- uno de cada por cada día de la
semana.
\end{solution}

\begin{solution}{exo:g1:day-and-night:9}
Podría tener razón: si nuestro lado siguiera mirando al Sol, aquí
seguiría siendo de día. Pero los niños del lado opuesto tendrían
entonces noche para siempre --- ¡no le darían las gracias a Zoe por su
idea!
\end{solution}
